\chapter{The bread next door}\label{connection:neighbours}
\regionline{Familiar ground, different work}
\atlasfiguretitle{The bread next door: paired comparison}
\comparisonplate{Focaccia}{Ftira}{A risen tray and a ring to fill. Their neighbours give familiar names different jobs.}

There is a particular pleasure in finding that a familiar name has become unfamiliar just down the coast. In Genoa, focaccia offers a risen dough dimpled with oil and brine. In Recco, it becomes two very thin sheets enclosing soft cheese. The same invitation to buy focaccia can lead to something held in a hand or something cut from a hot, fragile tray. Before asking which is the original, it is worth noticing what each preparation wants its baker to do.\citep{focacciaconsortium,reccofocaccia}

The Genoese dough needs time to grow. Its dimples keep the topping close to the surface, and its thickness leaves a crumb to tear. Recco's sheets need stretching rather than a rise. The cook watches their thinness and the cheese between them. Warmth is part of the experience because the filling and delicate crust belong together. These differences are large enough to change the baker's timetable, even while the breads share a name and a regional setting.

\section{A name can hold more than a loaf}
Malta offers a related pleasure. The ftira in the atlas is a flattened sourdough ring, baked before it is opened and filled. The Gozitan preparation discussed beside it can carry a topping into the oven. In one case, the bread waits for the meal to be assembled. In the other, the baking brings bread and topping together. A reader following only the word would miss the change in work.\citep{ftira,gozoftira}

This is a useful habit to bring to a bakery counter. Look for the local qualification, the filling, the surface and the way a portion is handed over. A food name may carry affection and familiarity without prescribing an exact formula. The varieties within a place are often the reason the name remains useful.

The arepa chapter makes that range especially inviting. Erica Dinho's Colombian examples include thin arepas paisas, sweet-maize arepas de choclo, cheese versions and arepas de huevo. Across the border, Irena Stein describes the Venezuelan arepa as a home for fillings and leftovers. These accounts offer particular makers and meals, rather than a border that neatly separates all thin breads from all thick ones.\citep{dinhoarepavar08,steinarepa08}

\section{A different day at the same table}
Baguette and pain de campagne let us compare two French loaves without turning difference into rivalry. A slender baguette gives a large amount of crust around its crumb. The country loaf printed here gives the baker a larger dough to ferment, fold and bake in a covered pot. Slice both and the changing balance of crust and interior is immediately visible. A crisp length broken for a meal and a substantial loaf cut for repeated slices offer different pleasures.

The stories also bring different kinds of continuity into view. UNESCO's baguette account follows skilled daily practice and the social life around bakeries. Poilâne's own history follows one bakery and its particular large loaf. Neither account needs to stand for all French bread to show how a loaf becomes expected in a neighbourhood. A customer learns where bread is made and when it will be ready.\citep{baguette,poilanestory,poilanemiche}

Rye gives that comparison another direction. In the Finnish account, eastern limppu appears as a round loaf while western reikäleipä is flatter and pierced for storage on poles. The name ruisleipä embraces both. Storage, shape and household rhythm sit inside the same bread family, so resemblance in grain tells only part of the story.\citep{finroot}

\section{Nearness is an invitation to look closer}
Aish baladi and Moroccan khobz are both useful companions to a meal, but their recipes in this book build different structures. Baladi is rolled with bran and puffed into a pocket. Khobz remains a low round whose crumb can be torn for dipping. The contrast is in what the finished bread lets a hand do, as well as in the dough that makes it possible. The distance between the two places is substantial, yet even the broad category of North African bread becomes richer when these practical differences are allowed to remain visible.\citep{saveurbaladi,benkhobz}

Neighbourhood matters because people share markets, ovens, names and meals. It does not remove the choices within them. A baker may be answering a need for a quick breakfast, a loaf to slice tomorrow or a bread that will carry its topping home. Following those choices gives an atlas something more generous than a row of national signatures. It lets one place contain several good answers.

\section{Different homes, shared heat}
The neighbourhood oven gives us a way to look beyond the differences on a counter. In the Fez account, households bring prepared dough on cloth-covered boards to a ferran. The workers recognise the boards and breads, watch the rise and return the baked loaves to their owners. Several households share the heat while keeping their own food recognisable.\citep{fezferran}

That arrangement makes a small distinction matter. A particular board, a shape or a remembered customer can tell the baker where a loaf belongs. The oven's work includes recognition, not simply loading dough. What appears to be one neighbourhood service can therefore hold the choices of many kitchens.

Azerbaijani tandir bread brings another oven-centred practice into view, while the qutab described beside it changes both the cooking and the meal. A filled half-moon can contain herbs, pumpkin or meat before it meets the heat. The substantial plain loaf waits to be torn beside other dishes. Both belong in the same wider food setting, yet the moment when bread and accompaniment meet is quite different.\citep{tandir,qutabvarieties}

For a baker reading an atlas, these accounts suggest questions to ask in any place: who mixes the dough, who owns the heat and how does a finished bread find its way back to the right table? Neighbouring differences can begin before the oven door opens. They can also survive the sharing of an oven, a market or a familiar bread name.

\begin{breadnote}{Bake two kinds of focaccia}
Follow the Genoese and Recco recipes on pages~25 and~26. Notice when most of the work happens, how the dough feels and what becomes of the tray at the table. Then read ftira on page~30 for another name whose local preparations change the relationship between bread and meal.
\end{breadnote}
